\section{Introduction}
\label{sec:introduction}
On small algorithmically generated tasks, fitting the training labels can precede held-out prediction by a long interval. This delayed generalization, or grokking, makes training accuracy an incomplete description of learning \citep{power2022grokking}. Prior experiments already vary transformer size and training duration on modular division and report model-wise grokking \citep{davies2023unifying}. We ask which size-dependent behavior and explanatory consequences hold in a precisely specified MLP system.

Our candidate explanation is that a wider trainable feature basis advances shared-structure learning while preserving a substantial interval after fit. Its consequences must be separated: earlier held-out onset, persistent post-fit delay, and a width benefit attributable to hidden trainability. A shorter generalization clock does not establish a shorter delay without the fit clock; a frozen model that cannot fit does not isolate feature learning from fitting capacity.

We test three prospective predictions on addition modulo 31, pairing fresh split/initialization blocks across five widths. All 204 registered runs include feature freezing, decay and learning-rate changes, unordered-orbit splits and separately trained multiplication models. We retain censored clocks, onset/terminal differences and seed-level uncertainty. The main comparison supports earlier restricted generalization and persistent delay, but the frozen fit gate fails (Section~\ref{sec:primary-results}). Complete controls expose recipe and threshold boundaries (Section~\ref{sec:scope}). The contribution is this controlled characterization and bounded explanatory test, rather than a first-size-study claim, a parameter-count law or a unique mechanism.
